\chapter{Identifying directional extensions} \label{ch:identifying2}

This chapter covers particular issues that arose in the process of identifying examples of the types of directional extensions presented in Chapter \ref{ch:identifying} and in creating a database that provides the basis for a quantitative approach to describing their distribution across Chadic languages. Some readers may prefer to first read Chapter \ref{ch:distribution} to see the quantitative results before reviewing the technical details provided in this chapter.

From a qualitative and descriptive perspective, this chapter provides insights into rarities found in directional extensions that are not captured by the quantitative generalizations in Chapter \ref{ch:distribution}. Section \ref{sec:multiple} covers languages where more than one morpheme expresses the same type of directional meaning. Conversely, Section \ref{sec:onetomany} describes cases where a single morpheme can express more than one directional meaning. Section \ref{sec:motionevent} explains what type of predications were considered motion events, specifically dealing with cases of transitive (caused) motion and fictive motion. From a methodological perspective, this chapter provides details about the sampling (Section \ref{sec:data}) and data processing behind the database (Section \ref{sec:database}). 

\section{Multiple forms for the same directional meaning} \label{sec:multiple}

There are a number of languages in which a particular directional meaning is expressed by more than one form in the verbal morphology. In many cases, the different forms are phonological alternations that can be analyzed as allomorphs of a single morpheme. These cases are uncontroversial and not discussed in this section. In most of the cases discussed below, the distribution of the different forms can be treated as a type of morphosyntactically conditioned allomorphy. However, there are also some instances where the different forms expressing the same directional meaning can be distinguished by an additional semantic factor. These are treated as separate directional extensions.

\subsection{Morphosyntactically conditioned allomorphy} \label{sec:allomorphy}

Of the 15 West Chadic languages with a ventive extension (Chapter \ref{ch:distribution}, Section \ref{sec:vent2}), ten have two forms which are distinguished by the tense-aspect of the clause (See Figure \ref{fig:westventive} in Chapter \ref{ch:diachrony}). In these languages, there is a form featuring the consonant \textit{n} which is said to be used in perfective contexts, and another form, frequently featuring \textit{t}, for non-perfective contexts \citep[307--309]{schuhChadicCornucopia2017}. Since tense-aspect is regularly expressed through verbal morphology in Chadic languages, these forms are treated in this book as a type of allomorphy and categorized as a single directional extension in each language. 

In \hyperref[mada]{Mada}, there are two forms of the ventive extension which are used for different verb classes. Most verbs are marked for ventive with the suffix \textit{-re} while a minor class uses palatalization to indicate ventive forms \citep{ernstkurdiTenseAspectMood2016}. These are treated as lexically-conditioned allomorphy, and therefore as realizations of the same morpheme.

There is a unique feature of transitive motion constructions found in four North Biu-Mandara languages.\footnote{Three of these languages are in the Mandaraic group of North Biu-Mandara languages. There is also some evidence suggesting that this may be the case in Hdi (North Biu-Mandara) as well. If so, this could help clarify the morphological analysis of Hdi which is in need of additional research before reliable cross-linguistic comparisons can be made.} An alternate form of the directional extension may be used when the figure on the path of motion is the object. For example, \hyperref[glav]{Glavda} has at least five directional extensions. Two of these have alternate forms depending on the transitivity of the verb. The upward suffix is \textit{-it} with intransitive predicates, as in (\ref{ex:glav18}), and \textit{-dit} with transitive predicates, as in (\ref{ex:glav7b}). 

\ea\label{ex:glav18}
\langinfo{Glavda verb \textit{ďal-} `climb' with \textit{-it}  suffix }{}{\cite[]{owensGlavdaRuralUrban2011}} \\
\gll ďal-it	\\
climb-\gl{up}		\\
\glt `climb up something (tree, mountain)'

\ex\label{ex:glav7b}
\langinfo{Glavda verb \textit{ɬəg-} `push' with \textit{-dit}  suffix }{}{\cite[663]{bubaGlavdaMorphology2007}} \\
\gll ɬəg-a-d-ít-ɬəga \\
push-\gl{3}-\gl{tr}-\gl{up}-push \\
\glt `He pushed something up/from east to west.'
\z

In a similar pattern, the Glavda downward suffix is \textit{-xi} with intransitive predicates and \textit{-di} with transitive predicates. There are also two forms given for the inward suffix. The suffix \textit{-dəm} only appears in two examples, both transitive verbs, while the suffix \textit{-m} appears with both intransitive and transitive verbs. In addition, the itive suffix \textit{-da} is only used with transitive predicates. 

Similar patterns are also found in \hyperref[dghw]{Dghwede}, \hyperref[psik]{Psikye} and \hyperref[park]{Parkwa}. \hyperref[dghw]{Dghwede} has as many as 11 directional extensions. Five of these extensions have an alternate form preceded by \textit{d}. As \citet[19]{frickVerbalSystemDghwede1978} notes, the forms beginning with \textit{d} are only used with transitive predicates, and often add a ventive directional meaning to the translation. There are, however, also examples of forms without the initial \textit{d} used in transitive motion contexts. 

\hyperref[park]{Parkwa} has at least seven directional extensions, with multiple forms attested for the ventive, upward, downward, inward and outward extensions.  \citet[89, 92]{jarvisEsquisseGrammaticalePodoko1989} sometimes parses the forms into multiple morphemes and sometimes treats them as single units, but there is no overall explanation given for the distribution of the different forms. However, it does appear that transitivity is a key factor, and it seems that another relevant factor is whether the object is definite or indefinite. 

\hyperref[psik]{Psikye} has as many as nine directional extensions. These include two upward extensions, \textit{-(a)te} and \textit{-(a)me}, and two downward extensions, \textit{-(i)yi} and \textit{-a(w)a}, and an across (or outward) extension \textit{-ve}. Each of these forms also appear preceded by \textit{-ak} with a slight exception for the downward extension \textit{-(i)yi} which is preceded by \textit{-a} becoming \textit{-ayi}. The forms with \textit{-ak} are all found in transitive constructions, and they express directed caused accompanied motion, rather than having a purely directional function (Chapter \ref{ch:CAM}).  %For example, the \hyperref[psik]{Psikye} suffix \textit{-akate} used to expressed upward subsequent caused associated motion is treated as a case of co-expression in regards to the morpheme \textit{-ate} which can express upward directional meaning. 

In the current study, all alternations between forms of directional extensions that appear to be determined by transitivity are treated as allomorphs of the same directional extension.

\subsection{Semantically overlapping directional extensions}

In other languages, multiple forms expressing the same directional meaning are treated as different morphemes. In \hyperref[kana]{Kanakuru}, there is a ventive suffix \textit{-tə(ru)}, as in (\ref{ex:kana1}), as well as a ``ventive equivalent auxiliary'' \textit{bo} which has a different morphosyntactic distribution, occurring before the verb, as in (\ref{ex:kana9}) \citep[76]{newmanKanakuruLanguage1974}. In light of their different morphosyntactic profiles, these are treated as separate morphemes in spite of their overlapping directional semantics.

\ea\label{ex:kana1}
\langinfo{Kanakuru verb \textit{ma} `return' with ventive suffix}{}{\cite[73]{newmanKanakuruLanguage1974}} \\
\gll à	ma-tə-ni \\
\gl{pfv} return-\gl{vent}-\gl{pfv} \\
\glt `He returned here.'

\ex\label{ex:kana9}
\langinfo{Kanakuru verb \textit{ma} `return' with \textit{bo} auxiliary}{}{\cite[76]{newmanKanakuruLanguage1974}} \\
\gll shìi	bo	ma-ma	\\
\gl{3}\gl{sg}	\gl{vent}	return-return	 \\
\glt `He is returning (coming back).'
\z 

In \hyperref[buwa]{Buwal}, there are two ventive extensions described as having a distance-related semantic distinction: \textit{-ā} `ventive proximal' and \textit{-xā} `ventive distal' \citep{viljoenGrammaticalDescriptionBuwal2013}. No examples are given of a proximal-distal distinction in a directional function, but this distinction is demonstrated for the subsequent motion interpretation of the ventive markers (Chapter \ref{ch:functions}, Section \ref{sec:covent}). In addition, the ventive markers are described as having different aspectual functions (Chapter \ref{ch:functions}, Section \ref{sec:tense}). Since the two forms have distinct semantic profiles, and it is not clear that they can be further parsed into multiple morphemes, they are treated (in this book) as two different extensions which can both express ventive direction. 

Two other languages have developed extensions with unique semantic components. In \hyperref[psik]{Psikye}, there are two upward extensions with slightly different meanings. The suffix \textit{-me} is used for steep upward direction (e.g., up a tree), while \textit{-ate} is used for less steep inclines (e.g., up a hill) \citep{smithKapsikiLanguage1969}. In \hyperref[mata]{Matal}, as discussed in Chapter \ref{ch:identifying}, Section \ref{sec:bounded1}, there are two inward extensions. The form \textit{ába} is used for uncovered containers and \textit{águ} is used for covered containers \citep[76]{dougopheGrammarMatal2021}. In these cases, there is a relevant semantic component that distinguishes the two forms. For this reason, they are treated in this book as separate extensions. Therefore, \hyperref[psik]{Psikye} has two upward extensions and \hyperref[mata]{Matal} has two inward extensions. 

\section{One-to-many form-function relationships} \label{sec:onetomany}

There are several senses in which a directional extension can be said to express more than one function. Section \ref{sec:portmanteau} covers directional extensions that simultaneously express more than one directional axis of a path of motion (synexpression). Section \ref{sec:ambiguous} covers directional extensions that express different directional meanings in different contexts (co-expression). This section is only concerned with the expression of directional meanings. Co-expression of other meanings with directional meaning are described in Chapter \ref{ch:functions}. 

\subsection{Synexpressive directional extensions} \label{sec:portmanteau}

Typically, a directional marker is described as expressing direction along a single directional axis. In contrast, \citet{frickVerbalSystemDghwede1978} describes complex directional suffixes in \hyperref[dghw]{Dghwede} which combine ventive or itive direction with another directional axis. For example, in a few examples of the suffix \textit{-àwilè}, this suffix is shown to have an outward ventive directional meaning, as in (\ref{ex:dghw9}). 

\ea\label{ex:dghw9}
\langinfo{\hyperref[dghw]{Dghwede} verb \textit{xwáy} `run' with  ventive outward suffix}{}{\cite[23]{frickVerbalSystemDghwede1978}} \\
\gll kâ' xwáy-àwilè	\\
\gl{narr} run-\gl{vent}.\gl{out} \\
\glt `He came running out.'
\z 

\citet{haspelmathCoexpressionSynexpressionPatterns2023} refers to the simultaneous expression of two distinct meanings in a single form as \textit{synexpression} (in analogy with co-expression). In an example of synexpression involving two directional meanings, the verbal extension is counted as expressing both meanings. Therefore, (\ref{ex:dghw9}) is an example of both a ventive extension and an outward extension. Note that Dghwede also has a more standard ventive marker whose only directional meaning is solely to indicate a path of motion toward the deictic center. 

There are other suffixes in Dghwede that are said to be synexpressive. The suffix \textit{-dáyà} is said to be a downward ventive marker and the suffix \textit{-tígè} is said to be a downward itive marker. However, when examples are given of these suffixes in context, the translations provided do not explicitly substantiate this description of their meaning.

\subsection{Co-expression of more than one direction} \label{sec:ambiguous}

As already mentioned in Chapter \ref{ch:introduction}, Section \ref{sec:comparing}, directional extensions are frequently co-expressive. A single form is found to express different meanings in different contexts. As will be discussed in detail in Chapter \ref{ch:functions}, co-expression typically involves one type of directional meaning with some other non-directional meaning (e.g., associated motion or tense-aspect). However, in a few cases, a verbal extension has one type of directional meaning in one context but expresses a different type of directional meaning in another context.

In one language, \hyperref[gude]{Gude}, the same morph \textit{-gərə} that expresses inward direction is also said to express downward direction, as in (\ref{ex:gude6}) and (\ref{ex:gude7}), and the morph \textit{-gi} that expresses outward direction also expresses upward direction, as in (\ref{ex:gude8}) and (\ref{ex:gude9}). \citet{hoskinsonGrammarDictionaryGude1983} does not discuss how these directional meanings are disambiguated, but the language also has a verb \textit{dəmə} `go in or out' in addition to potentially having further adverbial phrases to clarify the direction when necessary. For the purposes of this study, \textit{-gərə} is counted as an inward or downward directional extension according to the translation given, and \textit{-gi} is likewise counted as an outward or an upward directional extension as indicated by the translation.

\ea\label{ex:gude6}
\langinfo{\hyperref[gude]{Gude} verb \textit{dəmə} `enter/exit' with inward/downward suffix }{}{\cite[100]{hoskinsonGrammarDictionaryGude1983}} \\
\gll dəmə-gərə	\\
enter/exit-\gl{into}/\gl{down} \\	
\glt `go in'

\ex\label{ex:gude7}
\langinfo{\hyperref[gude]{Gude} verb \textit{dzə} `go' with inward/downward suffix }{}{\cite[25]{hoskinsonGrammarDictionaryGude1983}} \\
\gll dzə-gərə	\\
go-\gl{into}/\gl{down}	\\
\glt `go down'

\ex\label{ex:gude8}
\langinfo{\hyperref[gude]{Gude} verb \textit{dəmə} `enter/exit' with outward/upward suffix }{}{\cite[100]{hoskinsonGrammarDictionaryGude1983}} \\
\gll dəmə-gi \\
enter/exit-\gl{out}/\gl{up}	\\
\glt `go out'

\ex\label{ex:gude9}
\langinfo{\hyperref[gude]{Gude} verb \textit{ndərə} `climb' with outward/upward suffix }{}{\cite[25]{hoskinsonGrammarDictionaryGude1983}} \\
\gll ndərə-gi	\\
climb-\gl{out}/\gl{up} \\
\glt `climb up'
\z 

In a few languages, the outward extension is said to have an `across' interpretation, as in the \hyperref[marg]{Margi} example in (\ref{ex:marg20}) compared with (\ref{ex:marg6c}). Also, in \hyperref[psik]{Psikye}, the suffix \textit{-ve} is described as meaning `to go or to come across, to come up out of' \citep[118]{smithKapsikiLanguage1969}. These descriptions of co-expression of an `across' meaning with outward direction are not well attested in the available data and so are not considered as a separate type of directional meaning. 

\ea\label{ex:marg6c}
\langinfo{\hyperref[marg]{Margi} verb \textit{mə̀̀} `run' with the outward suffix \textit{-bá}}{}{\cite[122]{hoffmannGrammarMargiLanguage1963}} \\
\gll mə̀-bá	\\
run-\gl{out}		\\
\glt `to start up and run out'

\ex\label{ex:marg20}
\langinfo{\hyperref[marg]{Margi} verb \textit{ŋkə̀} `swim' with outward suffix}{}{\cite[123]{hoffmannGrammarMargiLanguage1963}} \\
\gll ŋkə̀-bá \\
swim-\gl{out} \\
\glt `to take p. across the river by swimming (also: to swim across ?)'
\z 

%\subsection{Co-expression with other functions} \label{sec:coexpression}

%A final note on identifying directional extensions is that the approach of this study is strictly synchronic. In a description of {\hyperref[maka]{Makary Kotoko}, \citet[239]{allisonGrammarMakaryKotoko2020} identifies a number of ``locative particles''. Allison proposes that the particles in (\ref{ex:maka1}) may be analyzed as follows: ``\textit{=he} conveys a downward direction of the action of the verb... and \textit{yo} that the action of the verb occurs moving away from the point of reference.'' However, the translations of the examples given do not include a translocative motion event. 
	
%	\ea\label{ex:maka1}
%	\langinfo{\hyperref[maka]{Makary Kotoko} verb \textit{gá} ‘put’ with ``locative particles''}{}{\cite[240]{allisonGrammarMakaryKotoko2020}} \\
%	\ea
%	\gll  gá =he \\
%	put =L.P.	\\
%	\glt `build'
%	
%	\ex 
%	\gll gá yo	\\
%	put L.P.	\\
%	\glt `diminish, reduce (in quantity)'
%	\z
%	\z 
	
%What Allison identifies in these examples are idiomatic verb-extensions combinations, and it is entirely possible that the extensions used were once directional extensions in Makary Kotoko. As discussed in Chapter \ref{ch:functions}, it is fairly common for directorial extensions to also have idiomatic uses. For example, in \hyperref[saku]{Sakun}, the verb root \textit{tʃíɗ} `leak' with the downward suffix becomes \textit{tʃíɗ-xá} `filter' and with the upward suffix becomes \textit{tʃíɗ-má} `coil' \citep[119]{thomasGrammarSakunSukur2014}. However, in Makary Kotoko there are no available examples which show the verbal extensions \textit{he} and \textit{yo} contributing information about direction to a predicate that describes a translocative motion event. For this reason, they are not counted (synchronically) as directional extensions.
	
\section{Identifying motion events} \label{sec:motionevent}
	
Directional extensions are defined as specifying the orientation of the path of motion. This definition entails that a morpheme can only have a directional function when it combines with a predicate describing a motion event. In this sense, a directional function contrasts with the function of associated motion\index{associated motion} \citep{guillaumeAssociatedMotionSouth2016}, as briefly discussed in Chapter \ref{ch:introduction}, Section \ref{sec:motion} and discussed in detail in Chapter \ref{ch:functions}. In principle, the distinction between directional meaning and associate motion is clear, but it depends crucially on diagnosing motion events. This section covers two situations where there may be disagreement about what to consider a motion predicate: transitive motion and fictive motion. 
	
In transitive (caused) motion events the figure on the path of motion is not necessarily the agent.\footnote{For a detailed discussion of the parameters of transitive directed motion, see \citet{osswaldDescriptionTransitiveDirected2022}.} For example, in (\ref{ex:cuvo0}), (\ref{ex:cuvo01}) and (\ref{ex:bura03}) the event described involves the subject causing the object to move.
	
	\ea\label{ex:cuvo0}
	\langinfo{\hyperref[cuvo]{Cuvok} verb \textit{kɛ̀ɮ} `throw' with ventive suffix}{}{\cite[303]{dadakGrammaireCuvokLangue2021}} \\
	\gll kɛ̀ɮ-ɛ́k	\\
	throw-\gl{vent}	\\
	\glt `Throw it here.' [French: `\textit{jeter vers ici}']
	
	\ex\label{ex:cuvo01}
	\langinfo{\hyperref[cuvo]{Cuvok} verb \textit{kə̀ɗ} `hit' with itive suffix}{}{\cite[311]{dadakGrammaireCuvokLangue2021}} \\
	\gll Kàdàmà á-kə̀ɗ-áɗ bàlɔ̀ŋ	\\
	Kadama \gl{3}\gl{sg}.\gl{sbj}-hit-\gl{itive} ball	\\
	\glt `Kadama sends the ball away by hitting it.' [French: `\textit{Kadama envoie le ballon loin en le tapant}.']
	
	\ex\label{ex:bura03}
	\langinfo{\hyperref[bura]{Bura} verb \textit{bəl} `break' with inward suffix}{}{\cite[8]{blenchBuraVerbalExtensions2010}} \\
	\gll bəl-wa	\\
	break-\gl{into}		\\
	\glt `to break articles into a receptacle'
	\z 
	
In a cross-linguistic study of directional and associated motion\index{associated motion} in verbal morphology, \citet[62]{rossCrosslinguisticSurveyAssociated2021} takes the position that caused motion events, as found in (\ref{ex:cuvo0}), (\ref{ex:cuvo01}) and (\ref{ex:bura03}), are not cases of directional meaning, but should be considered associated motion, specifically subsequent motion of the object. In this view, there is a distinction to be made between describing a translocative motion event and describing an object being ``set in motion''. On the other hand, in a study of directional and associated motion meaning in Tibeto-Burman languages, \citet[357]{genettiDirectionAssociatedMotion2020} assume a very broad view of motion events coding as motion verbs those with meanings such as `cut (away)', `saw (off)', and `tear (out)', and others similar to (\ref{ex:cuvo0}), (\ref{ex:cuvo01}) and (\ref{ex:bura03}).
	
The approach taken in this study differs slightly from the approaches of \citet{rossCrosslinguisticSurveyAssociated2021,rossPseudocoordinationSerialVerb2021} and \citet{genettiDirectionAssociatedMotion2020}. In contrast with Ross's definition, any example with a verb glossed to indicate that the object is necessarily set in motion (e.g., `throw', `send', `pour') is considered to describe a translocative motion event, in which case it is possible for a verbal extension to specify the direction of that path (as in \ref{ex:cuvo0}). In addition, a verb glossed to indicate that force is applied to an object such that it could potentially be set in motion (e.g., `hit', `push', `pull') could be used either to describe an event involving translocative motion (e.g., `hit a ball') or a non-motion event (e.g., `hit a wall'). Each example must be examined on a case-by-case basis to assess whether the event described can be assumed to describe translocative motion, irrespective of the presence of a directional extension. If the event described does include the description of a figure on a path of motion, it is possible for an extension to have a directional function, as in (\ref{ex:cuvo01}). 

In addition, it is assumed that any verb glossed to indicate that it describes an event of breaking or cutting does not entail that the object is a figure on a path of motion. In (\ref{ex:bura03}) the verbal extension adds a translocative motion event that would not be a part of the description without the verbal extension. Following the definition of \citet{guillaumeAssociatedMotionSouth2016}, if a morpheme adds the fact of motion to the utterance that would not otherwise be a part of the meaning, it is not considered to have a directional function. For this reason, in examples like (\ref{ex:bura03}) where an extension appears with a verb of cutting or breaking, any motion meaning in the translation is considered a resultative meaning, not a directional meaning.\footnote{This contrasts with the view of \citet{goldbergEnglishResultativeFamily2004} who group all directional meaning as part of their category of resultative. For the purpose of this study, resultative meaning is not considered a type of associated motion\index{associated motion} either. Associated motion typically involves a figure depicted as moving autonomously, whereas the resultative meaning is limited to the context of caused motion.} 

There is another context where there have been different perspectives on what counts as a directional function. In (\ref{ex:haus01}) and (\ref{ex:lagw0}) the translation contains an apparent directional meaning, but no argument of the verb can be identified as a figure on the path of motion. 
	
	\ea\label{ex:haus01}
	\langinfo{\hyperref[haus]{Hausa} verb \textit{har̃b} `shoot' with ventive suffix}{}{\cite[662]{newmanHausaLanguageEncyclopedic2000}} \\
	\gll bā sā̀ har̃b-ô-wā	\\
	?? ?? shoot-\gl{vent}-\gl{nmlz}		\\
	\glt `They are not shooting (at it) in this direction.'
	
	\ex\label{ex:lagw0}
	\langinfo{\hyperref[lagw]{Lagwan} verb \textit{ŋgù} `see' with inward suffix}{}{\cite[49]{lukasLogoneSpracheImZentralen1936}} \\
	\gll wā́-ŋgù-lí	\\
	??-see-\gl{into}\\
	\glt `look inside' [German: `\textit{sah hinein}']
	\z 

In (\ref{ex:haus01}), the event described includes a physical object moving through space, as entailed by the verb `shoot'. In the few cases where examples like these are found in the data on Chadic languages, the contribution of path meaning to that motion by a verbal extension is considered a directional function, even though the moving figure is not necessarily a syntactic argument of the verb. In contrast, in (\ref{ex:lagw0}), there is no identifiable figure on a path of motion. This metaphorical motion, or fictive motion \citep[115]{talmyCognitiveSemantics2000}, does not describe a literal figure in motion. The few cases of fictive motion in the data are not counted as directional (or associated motion).

\section{Data sources} \label{sec:data}

This study takes a maximal approach to sampling  information about the morphosyntax of Chadic languages. The primary goal is to be broadly descriptive, and only secondarily to examine evidence for any statistically significant patterns where possible. 91 of the 207 Chadic languages were included, as listed in the appendices and quantified in Table \ref{tab:dataset}.\footnote{In Table \ref{tab:dataset}, the circles in the column labeled \% are pie charts in which the black portion represents the percentage of total languages in each row that are included in the dataset. The numbers for Biu-Mandara North also count Jilbe (not included in this study), which Glottolog lists as an unclassified Biu-Mandara language.} These represent all of the languages for which grammatical descriptions of reliable quality with relevant information were accessible. For each subbranch there is data available for around half of the languages in the subbranch with the exception of West Chadic B where data and analyses are scarce.\footnote{For several West Chadic B3 languages there are reports which suggest that they do not have any directional verbal extensions. However, the available data were not considered substantial enough to accept that claim at face value. For that reason, those West Chadic B3 languages were excluded from the sample for this research.}

\begin{table}
	\caption{Number of Chadic languages included in study by subbranch}
	\label{tab:dataset}
	\begin{tabular}{llll} % add l for every additional column or remove as necessary
		\lsptoprule
		Branch &  Included &  Not included & \% \\
		\midrule
		West A            &  22 & 23 &  \bwpie{22}{23} \\
		West B            &  7 & 28 &  \bwpie{7}{28} \\
		Biu-Mandara Hurza &  2 & 0 &  \bwpie{2}{0} \\
		Biu-Mandara North &  27 & 24 &  \bwpie{27}{24} \\
		Biu-Mandara South &  13 & 15 &  \bwpie{13}{15} \\
		Masa North        &  2 &  4 &  \bwpie{2}{4} \\
		Masa South        &  2 & 2 &  \bwpie{2}{2} \\
		East A            &  5 & 9 &  \bwpie{5}{9}\\
		East B            &  11 & 11 &  \bwpie{11}{11}  \\
		\midrule
		Total       &             91 &   116 &  \bwpie{91}{116} \\
		\lspbottomrule
	\end{tabular}
\end{table}

The number of languages included is only slightly less than the number of Chadic languages for which the Glottolog records that either a (reference) grammar or a grammar sketch has been published, as shown in Table \ref{tab:grammar}. The scale in Table \ref{tab:grammar} refers to the ``description level'' categories used in the GlottoScope feature of the Glottolog \citep{hammarstromSimultaneousVisualizationLanguage2018}.\footnote{Thanks to Harald Hammarström for compiling the numbers in Table \ref{tab:grammar}.} While less than half (44\%) of all Chadic languages are included in this study, the coverage is very close to maximizing all of the existing published resources. The amount of data available for this study indicates that \citet[242]{frajzyngierChadic2012} were somewhat pessimistic in their estimate that only a quarter of Chadic languages (40/160) have been described.

\begin{table}
	\caption{Descriptive coverage of the grammars of Chadic languages}
	\label{tab:grammar}
	\begin{tabular}{lll} % add l for every additional column or remove as necessary
		\lsptoprule
		Languages &	\%	& Highest state of documentation \\
		\midrule
		47	&	23.03\%	&	grammar  \\
		46	&	22.54\%	&	grammar sketch \\
		5	&	2.45\%	&	dictionary \\
		22	&	10.78\%	&	specific feature \\
		6	&	2.94\%	&	phonology \\
		2	&	0.98\%	&	text \\
		65	&	31.86\%	&	wordlist \\
		3	&	1.47\%	&	comparative \\
		1	&	0.49\%	&	socling \\
		7	&	3.43\%	&	overview \\
		\lspbottomrule
	\end{tabular}
\end{table}

Taking Glottolog’s reference list as a general measurement of descriptive coverage, it is clear that more publications have focused on Hausa than any other language. Of the 3,097 references related to Chadic languages, 1,487 (45\%) are related to Hausa, as shown in Figure \ref{fig:grammars}.\footnote{Note that these numbers may include some repeated entries and some erroneously classified references. I assume that these discrepancies are evenly distributed.} The number of publications on Hausa is 167 times higher than the average of 8.9 publications per language in the rest of the language family. The other 83 West Chadic languages have received the least attention in the descriptive literature (5.5 publications per language). The coverage of East Chadic (10 publications per language), Masa (9.4 references per languages) and Biu-Mandara languages (11.7 publications per language) is roughly similar when taking into account the number of languages in each branch. 

\pgfplotstableread{ % data 
	Label     Other Grammars
	{East (36)} 354 9
	{Biu-Mandara (80)} 878 57
	{Masa (10)} 59 5
	{West w/o Hausa (83)} 437 22
	{Hausa} 1401 86
}\testdata

\begin{figure}
	\begin{tikzpicture}
		\begin{axis}[
			xbar stacked,   % Stacked horizontal bars
			xmin=0,         % Start x axis at 0
			ytick=data,     % Use as many tick labels as y coordinates
			width= 10cm,
			height= 5cm,
			legend style={at={(0.97,0.4)}},{anchor=south east},
			yticklabels from table={\testdata}{Label}  % Get the labels from the Label column of the \datatable
			]
			\addplot [fill=blue!60] table [x=Other, meta=Label,y expr=\coordindex] {\testdata};
			\addplot [fill=green!80] table [x=Grammars, meta=Label,y expr=\coordindex] {\testdata};   % "First" column against the data index
			\legend{Other,Grammars}
		\end{axis}
	\end{tikzpicture} %width=6cm,
	\caption{Grammars and other publications listed in Glottolog 5.2} \label{fig:grammars}
\end{figure}

Filtering for references that Glottolog classifies as a grammar (the most in-depth level of grammatical description) there are considerably more grammars of Biu-Mandara languages than the other branches.\footnote{The number of ``grammar'' level descriptive works in Glottolog is potentially exaggerated by some cases where a publication is erroneously identified as a grammar. For example, a review of a grammar is sometimes mistakenly identified as a grammar. The Glottolog bibliography also contains some duplicate records.} This research on directionals in Chadic languages reflects this skewed coverage with more information available about the grammars of Biu-Mandara languages compared to West Chadic or East Chadic languages.

Reliance on published descriptions of grammars naturally means that this study is a comparison of \textit{doculects} as defined by \citet[342]{goodLanguoidDoculectGlossonym2013}: ``a linguistic variety as it is documented in a given resource''. The analysis is limited to what information is published, and thereby leaves some gaps, including potentially omitting relevant details about linguistic variation (whether intra-speaker or dialectal). While these descriptions are based primarily on published sources, they are made more accessible by being compiled (in the appendices) in a manner that provides transparency and reproducibility \citep{berez-kroekerReproducibleResearchLinguistics2018}. In most cases, no change is made to the published analyses other than adapting the terminology to that used in this book. In other cases, the previously published data are re-analyzed. For example, a claim that a morpheme has a particular function is occasionally refuted based on the data provided, or additional functions of a morpheme are found in the data which are not mentioned in a previously published description. In a few cases, the morphological parsing in the original analysis is critiqued and an alternate analysis is proposed.

\section{Database} \label{sec:database}

The quantitative cross-linguistic results reported in the following chapters are all calculated from a database that was populated from published sources. One of the files that make up the database contains the entire collection of linguistic examples extracted for this research. There are a total of 3,591 example sentences (or phrases) collated into the database. An effort was made not to collect duplicate sentences, or too many redundant examples. Table \ref{tab:examples} shows the number of sentences or phrases per directional extension included in this study for each subbranch of Chadic languages.\footnote{Not all of the examples included in the database were ultimately relevant for this study. There are 265 examples from Hdi (Biu-Mandara North) which was not included due to issues with the analysis and presentation of the data. There are also some examples of multiverb constructions, a topic excluded due to insufficient data.} Some morphemes only appear in a single example. At the other extreme, one morpheme (the ventive in \hyperref[mina]{Mina}) is attested in 83 examples in the database. The mean is 17.8 examples per directional extension. The average numbers of example sentences per morpheme shown in Table \ref{tab:examples} indicate the same pattern of more substantial documentation of Biu-Mandara languages, and more sparse documentation of West Chadic B and Masa South.

\begin{table}
	\caption{Number of example sentences per morpheme by branch}
	\label{tab:examples}
	\begin{tabular}{ll} % add l for every additional column or remove as necessary
		\lsptoprule
		Subbranch &  Average examples \\ & per morpheme \\
		\midrule
		West A            &       17.2 \\
		West B            &        6.5 \\
		Biu-Mandara Hurza &       20.0 \\
		Biu-Mandara North &       16.9 \\
		Biu-Mandara South &       24.2 \\
		Masa North        &       10.8 \\
		Masa South        &        5.0 \\
		East A            &       11.2 \\
		East B            &       29.0 \\
		\lspbottomrule
	\end{tabular}
\end{table}

Each example sentence collected from available sources was analyzed for the function of all directional extensions in the example. The database file labeled ``functions'' has a column ``morphemeID'' that identifies the relevant directional extension using a unique label, as shown in the sample of data in Table \ref{tab:database}.\footnote{The actual database contains additional non-essential columns with notes about the form and function which were used to help guide the process of building the database.} The columns ``verb'' and ``verbgloss'' indicate which verb that morpheme was found modifying, and the column ``Function'' indicates what function the morpheme has in the example. Where the function is labeled ``Unclear'', that example was excluded from any statistical counts for directional functions (see Chapter \ref{ch:interpretation}). Therefore, if all examples of a putative directional extension are labeled as ``Unclear'' in the database, that morpheme is excluded from the statistics completely.
	
	\begin{table}
		\caption{Example data rows from ``functions'' file in database}
		\label{tab:database}
		\begin{tabular}{llllll} % add l for every additional column or remove as necessary
			\lsptoprule
			morphemeID  & Function & Path & verb & verbgloss & exampleID \\
			\midrule
			zulgVENT &  Unclear &  & zlá & take & zulg016 \\
			zulgVENT &  SubsCAM & VENT & gə̀s & catch & zulg030 \\
			zulgVENT &  SubsCAM & VENT & tə̀f & draw(water) & zulg031 \\
			zulgVENT &  SubsCAM & VENT & sə̀kə́m & buy & zulg023 \\
			zulgVENT &  Unclear &  & p & put & gemz004 \\
			zulgVENT &  Unclear &  & pél & untie & gemz012 \\
			buraOUT &  DIR & OUT & buka & push & bura126 \\
			buraOUT	& DIR & 	OUT	& nkya	& swim	& bura016 \\
			buraOUT	& DIR	& OUT	& ntsu	& remove	& bura017\\
			buraOUT	& DIR	& OUT	& pwara	& lead	& bura020\\
			buraOUT	& DIR	& OUT	& si	& come	& bura021; bura129\\
			sokoVENT	& DIR	& VENT	& ən	& arrive	& soko001\\
			sokoVENT	& DIR	& VENT	& ot	& come	& soko005; soko006 \\
			\lspbottomrule
		\end{tabular}
	\end{table}

For motion-related functions, the column ``Path'' records the direction of the path of motion. Each unique combination of morpheme, verb, function and (where applicable) path, requires a new row in the database. In the column ``exampleID'', a unique ID is given for all of the examples in the database where this particular combination of properties is attested. 

	
	\begin{table}[t]
		\caption{Coding for function of directional extensions}
		\label{tab:code}
		\begin{tabular}{ll} % add l for every additional column or remove as necessary
			\lsptoprule
			{} & Function   \\ %table header
			\midrule
			1 &	Direction (DIR)     \\		\midrule
			2 & Prior associated motion (PriorAM) \\
			3 & Concurrent associated motion (ConcAM) \\
			4 &  Subsequent associated motion (SubsAM) \\
			5 &  Subsequent caused accompanied motion (SubsCAM) \\
			6 & Resultative motion (RESULT) \\ \midrule
			7 & Locative (LOC)  	\\ 	\midrule
			8 & Tense-aspect (TAM) \\
			9 & Totality (TOTAL) \\
			10 & Additive (ADD) \\
			11 & Beneficiary (BEN) \\
			12 & Split in two (SPLIT) \\ \midrule
			13 & Lexicalized (Lex) \\
			14	& Caused accompanied motion (CAM) \\
			15 & Buy or sell (BUY) \\
			16  & Sell (SELL) \\
			17 & Fact of motion (MOTION) \\  \midrule
			18 & Fictive motion (FICTIVE)  \\ 
			19 & Unclear \\
			\lspbottomrule
		\end{tabular}
	\end{table}
	
	%For directional and associated motion functions, and additional column identifies the direction of the path of motion. These are mainly the eight directions described above: ventive (VENT), itive (ITIVE), upward (UP), downward (DOWN), inward (INTO), outward (OUT), onto (ONTO) and under (UNDER). There are also some less common directions, discussed further in Chapter \ref{ch:distribution}, Section \ref{sec:otherDIR} including homeward (HOME), bushward (BUSH) and sideward (SIDE). In the case of portmanteau directionals, the two directional components are written together separated by a period, as in VENT.OUT. 
	
	%Each row in the database of functions also identifies what verb (its form and gloss) the verb occurs with when it has a particular function, and the identifying labels for each example where that extensions combines with that verb to express the function identified in the row. This database serves as the source of the statistical data presented in the following chapters. 

\largerpage[-1]
The list of functions used to categorize examples in the ``Function'' column is shown in Table \ref{tab:code} in parentheses following the full label used for that function throughout this book. In some of the quantitative results shown, the labels are conflated into broader categories, as indicated by the horizontal lines. The functions in rows two through six are all treated as subtypes of associated motion\index{associated motion}. The functions in rows eight through 12 are all treated as subtypes of grammaticalized meaning. The functions in rows 13 through 17 are all treated as lexicalized uses of extensions.\footnote{An
    additional function ``ADVERB'' appears in the database to mark morphemes that are not considered extensions and thereby excluded from this study. Morphemes marked ADVERB in the database are not counted in any of the statistical tables presented.
    }


The entire database is archived and available open access on Zenodo \citep{lovestrandDirectionalsChadic2023}. The Zenodo deposit also contains the Python scripts that were used to extract and calculate the statistics that appear throughout this book. These are included in an effort to further increase the reproducibility of this research \citep[48]{winterStatisticsLinguistsIntroduction2019}. Maps shown in the following chapters were created using the lingtypology package for R \citep{morozLingtypologyEasyMapping2017}.
